\chapter{Competition Structure}

This chapter describes the general schedule of the competition.

\section{Training}

In order to guarantee safe and fair training conditions, the training sessions are divided into time slots.
The number of teams allowed on the track at the same time and the length of the slots will be announced on the website before the competition.
The commission might change the slots and the number of teams on the track without further notice.
In case of clear violations of training slots, the commission may issue penalties which will be subtracted from the final score of the respective teams.
In case of repetitive violations of slots or if team members endanger other teams or their equipment, the commission may expel single team members or whole teams from the competition.

% \section{Qualifying}

% There will be no qualifying in \yearofchallenge{}.

\section{Competition}
\subsection{Parc Fermé}

15 minutes ahead of the competition, all teams must hand in their vehicles at the “parc fermé”. 
No modifications of the vehicles must be made after this point.
Batteries must be separated from the system, the vehicle must be switched off.
All external tools must be removed from the vehicle, all wireless communication on board of the vehicles (Wi-Fi, Bluetooth, etc.) must be switched off or removed, except for the remote control communication.
The remote control must be placed next to the vehicle in switched off state.
When handing in the vehicle, the teams must make a definite statement to the commission in which events they would like to participate.
This is to ensure a smooth execution of the competition.

\subsection{Start Scheduling System}
\label{start_scheduling}

A traffic-light-like start scheduling system will signal the teams when to pick up their vehicle at the “parc fermé” and begin to prepare for starting.
The traffic light will show the following stages:

\begin{itemize}
	\item \textbf{1. Red - Standby:} \newline
	      No preparation is necessary.
		  The vehicle must be parked at the “parc fermé”.

	\item \textbf{2. Yellow - Preparation:} \newline
	      The team must pick up their vehicle at the “parc fermé” and prepare for the next event.
		  The team may switch the vehicles batteries and change the mode of operation as described in Section \ref{event_selection}.
		  No further modification of the vehicle is allowed at this point.
	      The team will receive a penalty if an active wifi connection or other external tool is required to start the vehicle.

	      The vehicle must be placed in the start box, located at the beginning of the track (cf. Section \ref{start_box}).
		  When ready, the team may signal the commission to start the event.

	      A total time limit of 5 minutes is given for preparation.
		  If the time limit is reached, the the event will proceed to the next stage automatically and the attempt may be canceled (cf. Section \ref{start_box}).

	\item \textbf{3. Green - Start:} \newline
	      The start box gate will open and the commission will start the timing of the event.
	      After each event, the vehicle must be returned to the “parc fermé” immediately.
		  Batteries must again be separated from the system, the vehicle must be switched off.
		  The remote control must be placed next to the vehicle in switched off state.
\end{itemize}

\subsection{Start Box}
\label{start_box}

The start box is separated from the track by physical barriers.
Up to two team members are allowed to prepare the start of the vehicle in the start box.
To the front of the start box is an openable gate, marked with a traffic sign and a matrix barcode (cf. Section \ref{fig_start_box_markings}).
An attempt starts with the opening of the gate.
The start box exit can be separated from the track by a solid white line. This line may be crossed to enter the track.
The gate of the start box remains open for 30 s. An attempt is canceled if:

\begin{itemize}
	\item The vehicle has not been placed ready to start in the box when the gate opens.
	\item The gate is forced open by a vehicle.
	\item The vehicle fails to leave the box while the gate is open.
	\item RC mode is activated inside the start box.
\end{itemize}

Penalties will not be applied until the vehicle passes the start line (i.e. collisions in the start box or driving outside of the right lane before passing the start line do not reduce the overall score).
